\documentclass[12pt]{article}
\usepackage[T1]{fontenc}
\usepackage[utf8]{inputenc}
\usepackage{amsmath,amssymb,amsthm}
\usepackage[margin=2.7cm]{geometry}
\usepackage{booktabs}
\usepackage{graphicx}
\usepackage{xcolor}
\usepackage{tikz}

\setlength{\parindent}{1.4em}
\setlength{\parskip}{0pt}

\newtheorem{theorem}{Theorem}
\newtheorem{definition}{Definition}
\newtheorem{corollary}{Corollary}
\newtheorem{remark}{Remark}

\begin{document}

% ===== PAGE 1 =====
\noindent{\LARGE\textbf{Why Two Bumble Filters Cut Your Pool by 80 Percent}}\\[0.3em]
\noindent{\small Hakvin Vosteen \hfill 2026}
\vspace{0.3cm}

\noindent{\small\textit{Combinatorial model. Filter pass-rates are illustrative midpoints.}}
\vspace{0.7cm}

\noindent You set two filters on a dating app. Distance, age range. They feel innocent. Each one alone removes maybe twenty percent of your pool. You assume the combined effect is forty percent. The combined effect is closer to sixty-five percent. With six filters you have less than ten percent of the original pool. With twelve, statistically zero matches. The math is multiplicative, not additive, and your intuition is wrong because nobody has trained your intuition on this.

\vspace{0.5cm}
\begin{center}
\colorbox{white}{\includegraphics[width=0.78\textwidth]{hook_fig.pdf}}
\end{center}

\vspace{0.3cm}
\medskip\noindent\textbf{Setup.}\quad Each filter is independent (or roughly so). The probability of surviving $n$ filters is $\prod_{i=1}^n p_i$. For typical pass-rates per filter $\bar{p} \approx 0.7$--$0.85$, the surviving fraction shrinks geometrically.

\newpage

% ===== PAGE 2 =====
\begin{definition}[Filter Pool Collapse Curve]
For a pool of size $N_0$ subjected to $n$ independent filters with pass-rates $p_1, \ldots, p_n$, the surviving pool size is
\[
N(n) \;=\; N_0 \cdot \prod_{i=1}^n p_i \;\approx\; N_0 \cdot \bar{p}^n,
\]
where $\bar{p}$ is the geometric mean of pass-rates. The collapse is exponential in $n$.
\end{definition}

\bigskip

\begin{theorem}[Curse of Dimensionality for Compatibility]
For two random vectors in $d$-dimensional compatibility space with marginal probability $p$ of agreement on each dimension, the joint probability of agreement is $p^d$. As $d$ grows, $p^d \to 0$ for any $p < 1$.
\end{theorem}

\begin{proof}
Independence gives $P(A_1 \cap \ldots \cap A_d) = \prod p_i$. For identical marginals $p$, this equals $p^d$. Logarithm is linear in $d$ with negative slope $\log p$, so the probability decays exponentially.
\end{proof}

\bigskip

\medskip\noindent\textbf{Plugging in dating-app numbers.}
\begin{itemize}
\setlength{\itemsep}{0pt}
\item Distance filter: $\bar{p} \approx 0.85$ (most users somewhat nearby)
\item Age range filter: $\bar{p} \approx 0.70$ (narrow ranges typical)
\item Two filters: $0.85 \cdot 0.70 \approx 0.60$ \quad(40\% pool gone)
\item Add height: $\cdot 0.60 \approx 0.36$ \quad(64\% gone)
\item Add education, smoking, kids, religion: $\cdot 0.5^4 \approx 0.022$ \quad(98\% gone)
\end{itemize}

\begin{corollary}
After roughly six filters, you are statistically guaranteed to interact with less than $5\%$ of your original pool. After ten, less than $1\%$. The "perfect partner" filter is not a filter, it is an empty set.
\end{corollary}

\newpage

% ===== PAGE 3 =====
\medskip\noindent\textbf{Why your intuition fails.}\quad Human reasoning about probability is approximately additive, not multiplicative, for small numbers of independent events. We expect "two filters at $20\%$ loss each" to give $40\%$ loss. The actual answer requires knowing that the geometric (multiplicative) reduction $0.8 \cdot 0.8 = 0.64$ gives $36\%$ loss, not $40\%$. The discrepancy seems small at $n=2$. At $n=10$ it is the difference between intuition saying "half remaining" and reality saying "ten percent remaining."

\bigskip

\medskip\noindent\textbf{The Bridge.}
\[
\boxed{\;\log N(n) \;=\; \log N_0 \;+\; n \cdot \log \bar{p}\;}
\]
The log of pool size is linear in filter count with negative slope. This is the Curse of Dimensionality applied to mate selection. Information-theoretically: each filter adds $-\log_2 \bar{p}$ bits of selectivity. After $n$ filters you have demanded $-n \log_2 \bar{p}$ bits of agreement, which generally exceeds the information actually shared between any two random people in the pool.

\bigskip

\begin{center}
\colorbox{white}{\includegraphics[width=0.74\textwidth]{curse.pdf}}
\end{center}

\bigskip

\medskip\noindent\textbf{Practical implication.}\quad The strategy of "set strict filters to find the right person" is mathematically self-defeating beyond a small number of dimensions. Empirical research on dating-app outcomes (Pew, Hinge internal reports, OkCupid blog circa 2014--2018) confirms that match rates collapse with filter count and that long-term relationship outcomes do not improve proportionally. The optimal filter count is small, typically $2$--$4$, and the additional filtering gains less than the lost pool size.

\newpage

% ===== PAGE 4 =====
\noindent\textbf{Empirical anchors.}

\medskip
\begin{tabular}{lll}
\toprule
Quantity & Value & Source \\
\midrule
Pew dating-app filter usage (2024) & $58\%$ use $3+$ filters & Pew Research \\
OkCupid match rate vs filter count (2017) & $\sim$ exponential decay & OkBlog (archived) \\
Bumble premium feature: ``Filters'' & up to $7$ simultaneous & Bumble app docs 2026 \\
Reported median Tinder match-rate (men) & $\approx 1\%$ of swipes & various community surveys \\
\bottomrule
\end{tabular}

\bigskip

\noindent\textbf{References}

\medskip
\noindent[1] Bellman, R. (1961). \textit{Adaptive Control Processes: A Guided Tour.} Princeton. (Original ``curse of dimensionality.'')

\medskip
\noindent[2] Rudder, C. (2014). \textit{Dataclysm.} Crown. (OkCupid behavioral analysis.)

\medskip
\noindent[3] Pew Research Center (2024). The state of online dating in the U.S.

\medskip
\noindent[4] Schwartz, B. (2004). \textit{The Paradox of Choice.} HarperCollins. (Choice-overload framing.)

\vspace{0.5cm}

\noindent\textbf{Notes on parameter estimates}

\medskip
\noindent The pass-rates $\bar{p} \approx 0.7$--$0.85$ are illustrative midpoints. Actual values depend on the user's pool composition: a user in a dense urban area has higher per-filter pass-rates than in a rural area. The exponential structure of the Filter Pool Collapse Curve holds in all cases; the slope $\log \bar{p}$ is what changes.

\medskip
\noindent What the model does not address: filter dependencies (height and weight are correlated, so adding both is less restrictive than independence assumes), platform-imposed minimum-pool thresholds (some apps refuse to show empty results), and active de-filtering by users when the pool collapses (people remove filters when matches dry up). Those factors soften the curve in practice but do not change its qualitative shape.

\end{document}
